\documentclass[11pt]{article}

% Change "review" to "final" to generate the final (sometimes called camera-ready) version.
% Change to "preprint" to generate a non-anonymous version with page numbers.
\usepackage[review]{acl}

% Standard package includes
\usepackage{times}
\usepackage{latexsym}
\usepackage{graphicx}
\usepackage{subcaption}
\usepackage{dblfloatfix}
\usepackage{hyperref}
\renewcommand{\dbltopfraction}{0.99}       % allow a top float to fill 99% of the page
\renewcommand{\dblfloatpagefraction}{0.99}

% preamble
\usepackage{cuted}
\usepackage{capt-of}
\usepackage{booktabs}
\usepackage{todonotes}
\usepackage{amsmath}
\usepackage{bm}

% For proper rendering and hyphenation of words containing Latin characters (including in bib files)
\usepackage[T1]{fontenc}
% For Vietnamese characters
% \usepackage[T5]{fontenc}
% See https://www.latex-project.org/help/documentation/encguide.pdf for other character sets

% This assumes your files are encoded as UTF8
\usepackage[utf8]{inputenc}

% This is not strictly necessary, and may be commented out,
% but it will improve the layout of the manuscript,
% and will typically save some space.
\usepackage{microtype}

% This is also not strictly necessary, and may be commented out.
% However, it will improve the aesthetics of text in
% the typewriter font.
\usepackage{inconsolata}

%Including images in your LaTeX document requires adding
%additional package(s)
\usepackage{graphicx}

% If the title and author information does not fit in the area allocated, uncomment the following
%
%\setlength\titlebox{<dim>}
%
% and set <dim> to something 5cm or larger.

\title{Orrery: Visual Exploration of Embedding Spaces and Sparse Autoencoder Features}

%% Orrery: a Platform for Visualization of Embedding and Sparse Autoencoders Feature Spaces 

%% Orrery: a Platform for Visual Exploration of Embedding and Sparse Autoencoders Feature Spaces 

%% Orrery: a Platform for Visual Exploration of Embedding Spaces and Sparse Autoencoders Features  

%A Platform for Embedding Visualisation and Mechanistic Interpretability

% Author information can be set in various styles:
% For several authors from the same institution:
% \author{Author 1 \and ... \and Author n \\
%         Address line \\ ... \\ Address line}
% if the names do not fit well on one line use
%         Author 1 \\ {\bf Author 2} \\ ... \\ {\bf Author n} \\
% For authors from different institutions:
% \author{Author 1 \\ Address line \\  ... \\ Address line
%         \And  ... \And
%         Author n \\ Address line \\ ... \\ Address line}
% To start a separate ``row'' of authors use \AND, as in
% \author{Author 1 \\ Address line \\  ... \\ Address line
%         \AND
%         Author 2 \\ Address line \\ ... \\ Address line \And
%         Author 3 \\ Address line \\ ... \\ Address line}


\author{
Giacomo De Luca, \\
Department of Enterprise Engineering\\
University of Rome Tor Vergata\\
Via del Politecnico 1, 00133, Rome, Italy \\
  \texttt{\footnotesize deluca@ing.uniroma2.it} \\}


%\author{
%  \textbf{First Author\textsuperscript{1}},
%  \textbf{Second Author\textsuperscript{1,2}},
%  \textbf{Third T. Author\textsuperscript{1}},
%  \textbf{Fourth Author\textsuperscript{1}},
%\\
%  \textbf{Fifth Author\textsuperscript{1,2}},
%  \textbf{Sixth Author\textsuperscript{1}},
%  \textbf{Seventh Author\textsuperscript{1}},
%  \textbf{Eighth Author \textsuperscript{1,2,3,4}},
%\\
%  \textbf{Ninth Author\textsuperscript{1}},
%  \textbf{Tenth Author\textsuperscript{1}},
%  \textbf{Eleventh E. Author\textsuperscript{1,2,3,4,5}},
%  \textbf{Twelfth Author\textsuperscript{1}},
%\\
%  \textbf{Thirteenth Author\textsuperscript{3}},
%  \textbf{Fourteenth F. Author\textsuperscript{2,4}},
%  \textbf{Fifteenth Author\textsuperscript{1}},
%  \textbf{Sixteenth Author\textsuperscript{1}},
%\\
%  \textbf{Seventeenth S. Author\textsuperscript{4,5}},
%  \textbf{Eighteenth Author\textsuperscript{3,4}},
%  \textbf{Nineteenth N. Author\textsuperscript{2,5}},
%  \textbf{Twentieth Author\textsuperscript{1}}
%\\
%\\
%  \textsuperscript{1}Affiliation 1,
%  \textsuperscript{2}Affiliation 2,
%  \textsuperscript{3}Affiliation 3,
%  \textsuperscript{4}Affiliation 4,
%  \textsuperscript{5}Affiliation 5
%\\
%  \small{
%    \textbf{Correspondence:} \href{mailto:email@domain}{email@domain}
%  }
%}

\begin{document}

\maketitle


\begin{strip}
  \centering
  \includegraphics[width=\textwidth]{media/geometry7.png}
  \captionof{figure}{A screenshot of Orrery displaying the nebula mode on the WordNet dataset. Each point represents a gloss, each colour a different extracted topic. In the table, the first semantic search results for "geometry", also highlighted in the scatterplot. On the left, the visualization controls panel. On the header, search and analytics panels can be accessed using the icons, and the Collections and SAE pages through the navigation bar.}
  \label{fig:teaser}
\end{strip}


\begin{abstract}

\textbf{Orrery} is an open-source platform for interactive 3-dimensional visualization of embedding spaces, with native support for Sparse Autoencoders (SAE), built for data exploration and interpretability research. Given any textual, image, or vector dataset, a single pipeline - embedding, projection, clustering, automatic topic labeling and, optionally, the collection of per-document SAE feature activations - turns it into a 3D constellation that can be searched visually, semantically, lexically, or by SAE feature names; probes can be trained, configured and scored directly in the interface. SAEs are automatically downloaded, stored, indexed and rendered as navigable constellations in which entering a prompt highlights the features it activates. Individual features can be inspected inside a dashboard and  injected in the model for causal steering. The star-cartography-inspired interface sustains 60 fps on 250k points, including WebGL animations and volumetric nebula shaders, on 8 GB of RAM. Orrery is released under the Apache 2.0 license at \url{https://github.com/Giacomo-De-Luca/embedding-orrery}.

%their decoder vectors

\end{abstract}


%\textbf{Orrery} is an open-source platform for interactive 3-dimensional visualization of embedding spaces, with native support for Sparse Autoencoders (SAE), built for data exploration and interpretability research. Given any textual, image, or vector dataset (local or from HuggingFace), a single pipeline - embedding, projection, clustering, automatic topic labeling and, optionally, the collection of per-document SAE feature activations - turns it into an explorable 3D constellation that can be searched visually, semantically, lexically (BM25), or by SAE feature names; linear, SVR and MLP probes can be trained, configured and scored directly in the interface. SAEs are automatically downloaded, stored, and indexed, and their decoder vectors are rendered as navigable constellations in which entering a prompt highlights the features it activates. Individual features can be inspected inside a dashboard and directly injected in the model for causal steering. The star-cartography-inspired interface sustains 60 fps on 250k points, including WebGL animations and volumetric nebula shaders, on 8 GB of RAM. Orrery is released under the Apache 2.0 license at [URL].


\section{Introduction}

\begin{figure*}[t]
    \centering
    \includegraphics[width=\textwidth]{media/poetry.png}
    \caption{The SAE dashboard. On the top features related to "poetry" are searched semantically. For the selected feature, it's possible to inspect: label, density, top and bottom logits, density histograms and similar features and top activating documents. On the right steered and baseline response with single vector steering on the question "What is your favorite job?"}
    \label{fig:mywidefigure}
\end{figure*}




Ever since Mikolov made meaningful textual embedding spaces widely available \cite{mikolov2013distributed}, multiple platforms have been developed to explore them visually. The Embedding Projector \cite{smilkov2016embeddingprojectorinteractivevisualization} let users explore embedding spaces as scatter plots in 3-dimension, reduced with PCA or t-SNE \cite{van2008visualizing}. Its accompanying paper named three tasks for visualizer to solve: (1) \emph{exploring local neighborhoods}, (2) \emph{viewing global geometry and finding clusters}, and (3) \emph{finding meaningful directions}, such as the famous man:woman\,::\,king:queen analogy \cite{mikolov-etal-2013-linguistic}. The Projector solved just the first task: it offered no way to cluster vectors or recover directions inside the platform.

After the popularization of contextualized embeddings \cite{devlin2019bertpretrainingdeepbidirectional, reimers-gurevych-2019-sentence}, new platforms emerged to visualize embedding spaces at scale \cite{wang2023wizmapscalableinteractivevisualization, nomic2022atlas, johnson2024latentscope, ren2025embeddingatlaslowfrictioninteractive}. Those support clustering and automated topic labeling, solving the second task proposed in the Projector's list, but revert to 2D, trading navigation for scalability. 

%such as WizMap \cite{wang2023wizmapscalableinteractivevisualization}, Nomic Atlas \cite{nomic2022atlas}, Embedding Atlas  \cite{ren2025embeddingatlaslowfrictioninteractive} and Latent scope\cite{johnson2024latentscope}.

% Later: "handle millions of data points" ... Unlike the Projector, all those projects reverted back to 2D, trading navigation for scalability. 

The third task, \textit{finding directions in the original space}, remains crucial in interpretability research, yet no visualization platform addresses it. Linear structures emerge unsupervised in LLMs: the separation between factual true and false can be visualized and recovered through linear probes \cite{Marks2023TheGO}; refusal is mediated by a single direction in the residual stream \cite{arditi2024refusallanguagemodelsmediated}. 
%Orrery allows training and configuring linear and nonlinear probes directly inside the platform §\ref{app:probes}. 

Additionally, since the Projector's release, new interpretability techniques built on vector spaces have emerged, most notably Sparse Autoencoders (SAE) \cite{elhage2022superposition, bricken2023monosemanticity}. SAE expands the dimensionality of the model's representation space to disentangle superposition \cite{elhage2022superposition} and recover interpretable features \cite{templeton2024scaling}. Sae features can be used in downstream tasks for steering, classification, evaluation
and retrieval \cite{jiang2025interpretableembeddingssparseautoencoders,
deng2026qwenscopeturningsparsefeatures}. 

%SAE features can be used in downstream tasks, such as classification, steering, clustering, evaluation and retrieval \cite{jiang2025interpretableembeddingssparseautoencoders, deng2026qwenscopeturningsparsefeatures}. 

%Since each row of the SAE decoder matrix is a vector in the original space, 
SAE features can be also visualized: Anthropic published a dashboard and 2D scatterplots to navigate SAE features in a UMAP-projected space as early as 2023 \cite{bricken2023monosemanticity, templeton2024scaling}. Open-source tooling, however, reproduced the per-feature dashboard while leaving the map behind (\S\ref{sec:related}).

Orrery aims to integrate embedding visualization with interpretability, solving all three of the Projector's tasks. To embedding-space exploration it brings interpretability tooling: corpora become searchable by SAE features, and hypotheses about linear structure can be tested by training probes directly in the interface. To interpretability it brings the capabilities of modern visualization platforms: SAE feature spaces are downloaded, projected, clustered, and labeled; they are also connected in the same interface to feature inspection and causal steering.

%% in order to facilitate interpretability and safety research, and to allow an immediate and interactive access to large textual corpus, such as the one from digital humanities. 

%Orrery aims to integrate embedding visualization with interpretability through its two enabled workflows. The first one allows the frictionless ingesting of local or Huggingface dataset turning them in visualizable collections 


%both to allow better navigation and 




%Orrery aims to bridge the gap between interpretability research and the current generation of embedding-visualization platform, solving all three of the Projector's intended use cases through two pipelines. 

%% here it should be explained what the first pipeline is for, allowing exploration of embedding spaces
%The first pipeline mirrors the Emedding/Nomic Atlas one: (1) embed the dataset, 2) compute projections, (3) cluster the vectors, and (4) label the clusters with cTF-IDF or LLM \cite{grootendorst2022bertopicneuraltopicmodeling}. Orrery also adds a fifth optional step: collecting per-document SAE activations, enabling SAE-feature search over the text. The processed collection is explored in 3D, where it can be colored, filtered, and searched semantically, lexically, or by SAE feature names. Linear and MLP probes can be trained and scored directly in the interface.

%% the second pipeline is for the exploration of feature spaces
%The second pipeline treats SAE features themselves as collections: features are visualized either by projecting the raw decoder vectors or by embedding their autointerpreter-generated labels \cite{bills2023language}. Orrery automatically downloads and store SAEs from HuggingFace and feature-labels from Neuronpedia \cite{neuronpedia}, then computes projections and topics for them. The resulting galaxy can be browsed visually, searched semantically over labels, or probed by entering a prompt and watching the activated features highlighted. Any interesting feature can be inspected in a dashboard inspired by Anthropic's and Neuronpedia \cite{neuronpedia} with: top logits, density histograms, similar features, top activating documents. Features can be injected using ActAdd \cite{turner2024steeringlanguagemodelsactivation} in a chat interface to causally steer the model. 



%Ever since Mikolov made meaningful textual embedding spaces widely available \cite{mikolov2013distributed}, multiple platforms have been developed to explore them visually. The Embedding Projector \cite{smilkov2016embeddingprojectorinteractivevisualization}, integrated into TensorBoard and released in 2016, let users explore embedding spaces as scatter plots in 3-dimensional spaces, reduced with PCA or t-SNE \cite{van2008visualizing}. Its accompanying paper named three tasks for visualizer to solve: (1) \emph{exploring local neighborhoods}, (2) \emph{viewing global geometry and finding clusters}, and (3) \emph{finding meaningful directions}, such as the famous man:woman\,::\,king:queen analogy \cite{mikolov-etal-2013-linguistic}. Regardless of the intent, the Projector solved just the first task: it offered no way to cluster vectors or recover directions inside the platform.

%After the explosion of LLMs and contextualized embeddings, new platforms emerged to visualize embedding spaces as scale, such as Wiz Map \cite{wang2023wizmapscalableinteractivevisualization}, Nomic Atlas \footnote{https://atlas.nomic.ai/map/owt}, Embedding Atlas \cite{ren2025embeddingatlaslowfrictioninteractive} and Latent scope \footnote{https://enjalot.github.io/latent-scope/}. Those platforms handle millions of data points and support clustering and automated topic labeling, solving the second task proposed in the Projector's list. Unlike the Projector, all those projects reverted back to 2D, trading navigation for scalability. 

%The third task, \textit{finding directions in the original space}, remains crucial in interpretability research, and has, to our knowledge, not been addressed by any of the current visualization platform. Linear structures emerge unsupervised in LLM: the separation between factually true and false can be clearly visualized and recovered through linear probes \cite{Marks2023TheGO}; refusal is mediated by a single direction, recoverable by difference-in-means \cite{NEURIPS2023_d066d21c}  that can be ablated or added to remove or induce refusal behavior \cite{arditi2024refusallanguagemodelsmediated}.

%Additionally, since the Projector's release, new interpretability techniques built on vector spaces have emerged, most notably Sparse Autoencoders (SAE) \cite{elhage2022superposition, bricken2023monosemanticity}, which expand the dimensionality of the model's representation space to disentangle superposition \cite{elhage2022superposition} and recover interpretable features \cite{templeton2024scaling}. SAE features can be used in downstream tasks, such as classification, steering, clustering, evaluation and retrieval \cite{jiang2025interpretableembeddingssparseautoencoders, deng2026qwenscopeturningsparsefeatures}. 

%Since each row of the SAE decoder matrix is a vector in the original space, SAE features can be visualized: Anthropic published a dashboard and 2D scatterplots to navigate SAE features in a UMAP-projected space already in 2023 \cite{bricken2023monosemanticity, templeton2024scaling}. Open-source tooling, however, reproduced the per-feature dashboard while leaving the map behind (\S\ref{sec:related}).


%Orrery aims to bridge the gap between interpretability research and the current generation of embedding-visualization platform, solving all three of the Projector's intended use cases through two pipelines. 

%% here it should be explained what the first pipeline is for, allowing exploration of embedding spaces
%The first pipeline mirrors the Emedding/Nomic Atlas one: (1) embed the dataset, 2) compute projections, (3) cluster the vectors, and (4) label the clusters with cTF-IDF or LLM \cite{grootendorst2022bertopicneuraltopicmodeling}. Orrery also adds a fifth optional step: collecting per-document SAE activations, enabling SAE-feature search over the text. The processed collection is explored in 3D, where it can be colored, filtered, and searched semantically, lexically, or by SAE feature names. Linear and MLP probes can be trained and scored directly in the interface.

%% the second pipeline is for the exploration of feature spaces
%The second pipeline treats SAE features themselves as collections: features are visualized either by projecting the raw decoder vectors or by embedding their autointerpreter-generated labels \cite{bills2023language}. Orrery automatically downloads and store SAEs from HuggingFace and feature-labels from Neuronpedia \cite{neuronpedia}, then computes projections and topics for them. The resulting galaxy can be browsed visually, searched semantically over labels, or probed by entering a prompt and watching the activated features highlighted. Any interesting feature can be inspected in a dashboard inspired by Anthropic's and Neuronpedia \cite{neuronpedia} with: top logits, density histograms, similar features, top activating documents. Features can be injected using ActAdd \cite{turner2024steeringlanguagemodelsactivation} in a chat interface to causally steer the model. 

% a closing statement on who is orrery for. Interpretability researchers, or people visualising corpus such as the digital humanities ones. 

%TODO check if to keep contributions: 
%Contributions: 

%\begin{itemize}
%    \item An embedding collection manager with a double database architecture that stores datasets, projections, embeddings, metadata, topics and per-document SAE activations, and automatically downloads SAEs from Gemma-Scope-2 together with their Neuronpedia labels and activation data.
%    \item A star-cartography-inspired 3D interface, with custom WebGL volumetric shaders, animations, 3D labels, advanced filtering,  and four search modes: visual, semantic, lexical, and SAE-based.
%    \item An SAE workflow that unifies exploration and intervention: feature spaces browsable, clusterable, and labelable like any other embedding collection, with per-feature dashboards and in-interface causal steering via activation addition.
    %\item Two replication studies to show how the platform can enable research: using the methodology of
%\end{itemize}


\section{Related Works}
\label{sec:related}

\paragraph{Embedding-visualization platforms.}



The generation of platforms that followed the Embedding Projector (\S1) converged on two design choices: 2D projections, and scale as the primary goal. Nomic Atlas \cite{nomic2022atlas} renders millions of embeddings with clustering, topic labels, filtering and search, but is a closed commercial service with a fixed embedding model and projection model. WizMap \cite{wang2023wizmapscalableinteractivevisualization} and Embedding Atlas \cite{ren2025embeddingatlaslowfrictioninteractive} reach similar scale in the browser, adding density contours, cross-filtering, and analytical charts. Latent Scope, \cite{johnson2024latentscope} an
open-source local pipeline is the only platform with any SAE integration: documents can be filtered by features from an SAE trained on \texttt{nomic-embed-text-v1.5} embeddings \footnote{\url{https://enjalot.github.io/latent-taxonomy}}, but the feature space itself is never visualized. 

Three-dimensional embedding views survived mainly as demonstrations: Google's \textit{t-SNE Map} of
artworks\footnote{\url{https://artsexperiments.withgoogle.com/tsnemap}} and the Hugging Face community's \textit{Semantic Galaxy},\footnote{\url{https://huggingface.co/spaces/webml-community/semantic-galaxy}}
an in-browser 3D view of EmbeddingGemma \cite{vera2025embeddinggemmapowerfullightweighttext} vectors and the main inspiration for Orrery's cosmic theme. The third dimension is not cosmetic: \citet{9222098} show that rotation-based navigation, as opposed to pure zooming, increases immersion and supports cluster understanding.   

%None of these platforms supports training probes or recovering directions.

%They have also limited configuration setup: Nomic support only their closed-source embedding model and projection, Embedding Atlas only supports density clustering and cTF-IDF labels. 

%targets embedding-model rather than LLM representations,
%\paragraph{3D views.} Three-dimensional embedding views survived mainly as demonstrations: Google's \textit{t-SNE Map} of artworks \footnote{\url{https://artsexperiments.withgoogle.com/tsnemap}} and the Hugging Face community's \textit{Semantic Galaxy} 
%\footnote{\url{https://huggingface.co/spaces/webml-community/semantic-galaxy}}, an in-browser 3D view of EmbeddingGemma vectors and the main
%inspiration for Orrery's cosmic theme. The third dimension is not cosmetic:
%\citet{9222098} show that rotation-based navigation, as opposed to pure zooming, increases immersion and supports cluster understanding. %\todo{add a line here about the fact that orrery aims at }

%Both are fixed-purpose demos rather than platforms. 

\paragraph{SAE feature visualization.}


%% TODO add the paper from 2019 activation atlas and maybe https://www.goodfire.ai/research/evee-explaining-genetic-variants#

Activation atlases \cite{carter2019activation} mapped the activation space of
vision models on a projected grid. For SAEs, \textit{Towards Monosemanticity} \cite{bricken2023monosemanticity} established the visual vocabulary still in use: a per-feature dashboard (activation density, top and bottom logits, activating examples) accompanied by a 2D UMAP map of 4,608 features, later extended to Claude Sonnet \cite{templeton2024scaling}. Both interfaces are on single layers and closed-source. Open-source tooling, such as SAE-Vis\cite{sae_vis}, SAEDashboard\cite{sae_dashboard}, and Neuronpedia \cite{neuronpedia} replicated the dashboard leaving the navigable space behind. Neuronpedia offers search over feature explanations, activation testing, and in-browser steering, but its only spatial view is a 2D UMAP of GPT-2 Small, without clustering or navigation aids. Google's PAIR published a pedagogical explorable of a 2D visualization of a single Gemma Scope \cite{lieberum2024gemmascopeopensparse} SAE\footnote{\url{https://pair.withgoogle.com/explorables/sae}}. Existing feature maps are one-off showcases of a fixed SAE rather than research platforms. 




\section{System Design}
\label{sec:system}

\begin{figure}[t]
    \centering
    \includegraphics[width=\columnwidth]{media/pipeline4.png}
    \caption{pipeline}
    \label{fig:pipeline}
\end{figure}


Orrery's design goal is seamless exploration of embedding spaces: frictionless~\cite{ren2025embeddingatlaslowfrictioninteractive} movement from a raw dataset to a navigable constellation, and from a promising SAE feature to its causal effect on the model~\cite{templeton2024scaling}. The platform allows two workflows over one shared backend. The \emph{collection workflow} (embed, project, cluster, label, explore) turns any textual, image, or vector dataset into a searchable space. The \emph{SAE workflow} (visualize SAE feature spaces, examine in the dashboard and causal steering) reuses the same engine for interpretability. 

%A  resuing the same enging  (\S\ref{sec:sae}): visualize SAE feature spaces, examine in the dashboard and direct causal steering. 

%% cut: allowing discovery in the scatterplot,

\paragraph{Datasets, collections, storage.}
Orrery separates two concepts. A \emph{dataset} is a set of documents (text, images, or pre-computed vectors) stored and indexed in DuckDB. A \emph{collection} is one representation of that dataset: an embedding, its 2D/3D projections, optional topic extractions and SAE activations. One dataset can therefore have many collections: the same corpus embedded with different models, prompts, or column combinations. DuckDB acts as the central orchestrator (documents, metadata, projections, SAE activations), while ChromaDB \cite{chroma2023} stores only IDs and dense vectors for approximate nearest-neighbour (ANN) search. A GraphQL API serves the Next.js frontend; long-running jobs stream progress over WebSockets and are resumable after interruption.

% Document duplication was the main limitation encountered with the initial ChromaDB-only storage. 

\paragraph{Ingestion.}
A dataset is previewed from HuggingFace or from local JSON, CSV, or Parquet files. The columns to embed are specified with a free-form template that supports both custom prompts and column concatenation (e.g.\ \texttt{\{word\}: \{definition\}} for WordNet). Eight embedding providers are supported \ref{app:embedding}. All features run fully offline with local models. \emph{Projections}, the only mandatory step to visualize a collection, are computed with PCA and UMAP~\cite{McInnes2018}, in both 2D and 3D. 

\paragraph{Topic modeling.}

Topic extraction follows and forks BERTopic~\cite{grootendorst2022bertopicneuraltopicmodeling}: since embedding and reduction already exist in the platform, we reuse only its clustering, c-TF-IDF, and labeling components, including its LLM labeling prompt. Clustering (HDBSCAN, k-means, Gaussian mixtures, or spectral) runs on the stored 2D/3D projections, on 5-D UMAP as in BERTopic, or on the original vectors; with BERTopic's defaults, silhouette and $C_v$ match to the fourth decimal. Labels come from c-TF-IDF keywords or LLMs and can be regenerated or manually edited.

%Topic extraction follows BERTopic~\cite{grootendorst2022bertopicneuraltopicmodeling}: embed, reduce, cluster, extract keywords, label. Since embedding and reduction already exist in the platform, we do not import BERTopic itself but fork only its clustering, c-TF-IDF, and labeling components, reusing its LLM labeling prompt. Clustering (HDBSCAN, k-means, Gaussian mixtures, or spectral) can be run on the stored 2D/3D projections,  5-D UMAP as in BERTopic, or on the original vectors. With BERTopic's default settings the results are identical over silhouette and $C_v$ to the fourth decimal.  Labels come from c-TF-IDF keywords or LLMs and can be regenerated or manually edited. 

%Topic evaluation scores can be computed directly in the collection page.





%% TODO add in appendix the table with results

%An optional reduction step merges topics hierarchically toward a target count. The originals are counted as subtopics for nested colouring.

% Run with BERTopic's default recipe (HDBSCAN over a 5-dimensional UMAP), our implementation reproduces its silhouette and $C_v$ scores to the fourth decimal.\todo{Commit the parity-check artifact backing this number.}

\paragraph{Per-document SAE activations.}
Any stored collection can optionally be run through Gemma Scope 2 SAEs~\cite{mcdougall2025gemmascope2}. Following \citet{jiang2025interpretableembeddingssparseautoencoders}, token-level activations are max-pooled into a single sparse vector per document. Because each dimension carries an auto-interpreted feature label, the corpus becomes searchable by feature names. 

\paragraph{Search and filtering.}
A collection can be searched four ways, with results rendered as a highlighted constellation centered on the first result and coloured by similarity: (i)~\textbf{semantic} search over the original embedding space, also triggered by clicking any point; (ii)~\textbf{text} search (exact, partial, or BM25, with column selection); (iii)~\textbf{SAE feature search}: querying a SAE feature name (e.g.\ \emph{football}) ranks documents by their max activation over matching features; (iv)~\textbf{prompt highlighting}, on SAE collections (\S\ref{sec:sae}): the prompt is run through the model, encoded by the SAE at the selected layer, and every activated feature lights up, ranked and coloured by activation strength. This inverts the SAE search: there, features index documents; here, a document indexes features.


%via FTS extension
%executed server-side in DuckDB
%, - the camera flies to it and its nearest neighbours in the pre-projection space light up

%% TODO check this better. put ref to image.



%The legend doubles as a filter: categorical fields render as toggleable category lists with counts, numeric fields as histograms with draggable range handles, and auto-detected temporal fields get a dedicated range slider that mutes out-of-range points.


%%% TODO ref to an image. 
%\paragraph{Colour.}
%Beyond the standard d3 scales, Orrery ships Crameri's perceptually uniform scientific colour maps~\cite{crameri2018scientific}, which avoid scale distortion and remain readable for colour-blind users~\cite{crameri2020misuse}. The numeric legend is an interactive histogram of the field's distribution with draggable minimum, center, and maximum handles plus a log option, so outliers skewing a colour scale can be visually spotted and directly re-anchored. 
%a scale skewed by outliers is re-anchored in seconds. 


%% TODO understand if I want to keep this in, probably now. 
%For datasets that are themselves about colour, a new \emph{Hilbert} scale samples 512 points from an $8^3$ grid of RGB space and orders them along a 3D Hilbert curve, linearising colour space onto a single locality-preserving strip (used in \S\ref{sec:xkcd}). The active colour configuration round-trips through the URL and can be saved per collection.

%%% TODO link to the appendix.
\paragraph{Probing.}
Probes can be trained and configured directly in the interface on the original embeddings: any categorical metadata field serves as target, fitted with a linear probe (Mass-means, logreg, ridge), SVR or a small MLP.  Linear probes double as directions: the learned weight vector is stored with the collection, and points can be coloured by their projection onto it (\S\ref{sec:glasgow}).

% turning a hypothesis about linear structure into an inspectable gradient

\begin{figure}[t]
    \centering
    \includegraphics[width=\columnwidth]{media/whitesae.png}
    \caption{Navigable map of the embeddings of the Neuronpedia labels of the Gemma-scope-2 layer 29 SAE of Gemma-3-4B-it. Topic labels appear dinamically when zooming or rotating. Clicking on a point triggers a zoom in animation, similarity search and allows inspection in the dashboard.}
    \label{fig:saecluster}
\end{figure}


\paragraph{Rendering 3D galaxies at scale}


The scatter plots are rendered through WebGL using a fork of Plotly.js: on every trace update, Plotly re-reads all $n$ points across all $m$ attributes of every existing trace, an $O(n \cdot m)$ burden that makes real-time overlays and filters impossible at scale. Plotly is used only for camera control and point rendering: tooltips, legend, labels, and animations are re-implemented on an overlay canvas that read the WebGL context directly. Legends double as controls: beyond d3's scales, Orrery ships Crameri's perceptually uniform colour maps~\cite{crameri2018scientific, crameri2020misuse}, and numeric legends render as interactive histograms of the field's distribution, with draggable handles to re-anchor scales skewed by outliers \ref{fig:concreteness}.  Label placement ports the Embedding Projector's collision-avoidance code~\cite{smilkov2016embeddingprojectorinteractivevisualization}: 3D anchors are projected to screen space each frame and packed on a spatial grid, extended with a priority order and rank-dependent opacity. Since no scatter-plot library offers volumetric rendering, the \emph{nebula} mode approximates it: we render feathered sprites around cluster centroids, spread according to cluster extent, in the style of GalaxyThreeJS,\footnote{\url{https://github.com/pickles976/GalaxyThreeJS}}. 
For 2D density contouring we port the code from Embedding Atlas \cite{ren2025embeddingatlaslowfrictioninteractive}, adding coloring according to the categories. 


%greedily
%Tooltips are positioned from projected coordinates, and \textit{fly-to} animations manipulate the camera directly.

%after an unsuccessful attempt with three.js bloom, 

%(cluster labels first, then similarity-ranked point labels)
%composited over the plot in screen blend mode
%and replacing the real-time automatically generated topics with the high quality ones cached from the topic modeling step. 

%\todo{ ADD BENCHMARK measumements in the APPENDIX The interface sustains 60\,fps at 250k points with nebulae enabled on 8\,GB of RAM.\todo{State test machine and measurement method.}


\subsection{SAE collections}
\label{sec:sae}


%Since the SAE features can be mapped to vectors in the original vector space of the model by extracting the w\_dec row of the matrix, the internal space of the features is already a dataset that can be visualised and labeled using one of the autointerpreters \cite{paulo2025automaticallyinterpretingmillionsfeatures}. Alternatively, the labels of the autointerpreters can be directly embedded with a normal embedding model, offering a navigable visualization of the feature space. 

%TODO probably delete, repeated 
%%While this procedure was done in 2D by Anthropic, to our knowledge no open-source project has navigable scatterplots for SAE features. 

%When we visualise a galaxy for the SAE features we can employ a different search modality, which, given a prompt, highlights in the visualisation all the features that gets activated by the prompt, ranking and colouring by activation level. 

%Once an interesting feature is found, like the one for poetry, right click on it allows inspection in the \textit{SAE dashboard}. The feature can then be injected directly inside the model using ActAdd \cite{turner2024steeringlanguagemodelsactivation}. The discovery, inspection and intervention workflow can be completed entirely in the same interface. 

%From the same platform can be completed the discovery, inspection and intervention workflow. 

%%TODO probably delete

%which features the top and bottom logit, density distribution histogram, and the top activating documents. 

%\todo{add n_dimensional} This is so repeated. 
\paragraph{Feature spaces as collections.}

Each SAE feature is a row of the decoder matrix $W_{\mathrm{dec}}$ — a vector in the model's original space — so feature sets can be visualized like any vector dataset, by projecting the decoder vectors directly or by embedding their auto-interpreter labels \cite{bills2023language, paulo2025automaticallyinterpretingmillionsfeatures} with any provider from \S\ref{sec:system}. From the SAE tab in the \textit{Collections} page, the user selects the target model, layer, and width; the SAE is downloaded and registered as a dataset together with its feature labels and activation examples from Neuronpedia \cite{neuronpedia}. Since these are ordinary vector collections with metadata, the entire collection pipeline (Figure~\ref{fig:pipeline}) is reused with no additional code.
 

%The demo currently supports all the Gemma-Scope-2 SAEs \cite{mcdougall2025gemmascope2}, with upcoming support for the Qwen-Scope SAEs \cite{deng2026qwenscopeturningsparsefeatures}. 

%\paragraph{Prompt highlighting.}
%SAE collections enable a fifth search mode: given a prompt, the backend runs a forward pass, encodes the selected layer with the SAE, and lights up every activated feature in the constellation, ranked and coloured by activation strength. This inverts the corpus workflow's feature search: there, features index documents; here, a document indexes features.

\paragraph{Feature dashboard and steering.}
Any feature in the SAE galaxy can be right clicked to examine it into the SAE page: a dashboard inspired by \citet{bricken2023monosemanticity},  and Neuronpedia \cite{neuronpedia}. For each feature, the dashboard displays top and bottom logits, activation-density histograms, features with similar activation patterns, top or decile based activation samples. Features are located through textual or semantic search, or by prompting the model and examining token-level activations in a heatmap. The found features can be directly injected into the model with ActAdd \cite{turner2024steeringlanguagemodelsactivation}, selecting the activation coefficient. Steering with multiple features is supported, as is split view to compare with baseline response. 

%The activated features density-filtered to avoid always activating ones. There is also an optional coverage filter, removing features that appear on all the input tokens above a configurable threshold (default > 80\%).

\paragraph{Interpretability backend.}

SAELens/TransformerLens \cite{bloom2024saetrainingcodebase, nanda2022transformerlens} support neither multimodal inputs or streaming generation with interventions active. The requirements for SAEs are small, - read/write caches at three sites per layer (post-attention, post-MLP, residual stream), and a single layer MLP for each SAE, - so we implement it directly. \texttt{Gemma-pytorch} was forked to add activation caches and streaming generation. JumpReLU \cite{rajamanoharan2024jumpingaheadimprovingreconstruction} for Gemma-Scope and 
TopK \cite{gao2024scalingevaluatingsparseautoencoders} SAEs for Qwen-Scope \cite{deng2026qwenscopeturningsparsefeatures} are written in plain PyTorch. Qwen-Scope SAEs are run via \textit{register\_forward\_hook} on the Transformers model.

%as in the official documentation \footnote{\url{https://huggingface.co/Qwen/SAE-Res-Qwen3.5-27B-W80K-L0\_50}. 

%Activation collection and steering impose requirements the standard
%SAELens/TransformerLens stack \cite{bloom2024saelens,
%bloom2024saetrainingcodebase} does not meet: multimodal inputs (Gemma~3 accepts images), model families outside its support matrix, and streaming generation with interventions active, so steering works as a real-time chat. Since the required functionality is small, just read/write caches at three sites per layer (post-attention, post-MLP, residual stream), and single layer MLP for each SAE, we implement it directly\footnote{Hook libraries over stock HF models
%\cite{fiottokaufman2024nnsight, wu-etal-2024-pyvene} could likely meet these requirements and could be considered for future versions of Orrery}: a fork of the reference \texttt{gemma-pytorch}%
%\footnote{\url{https://github.com/google/gemma_pytorch}} with caches
%and streaming added, plus JumpReLU \cite{rajamanoharan2024jumpingaheadimprovingreconstruction} and
%TopK \cite{gao2024scalingevaluatingsparseautoencoders} SAEs in 135 lines of PyTorch. The interface
%exposes Gemma Scope 2 \cite{mcdougall2025gemmascope2} end-to-end;
%Qwen-Scope \cite{deng2026qwenscopeturningsparsefeatures} runs at the
%backend level only (see Limitations).

\section{Case Studies}
\label{sec:cases}




\begin{figure}[t]
    \centering
    \includegraphics[width=\columnwidth]{media/concretenesswhite.png}
    \caption{Concreteness}
    \label{fig:concreteness}
\end{figure}

%\begin{equation}
%p_{\mathrm{mm}}(x) = \sigma\!\big((\mu^{+}-\mu^{-})^{\top}x\big),
%\end{equation}

We validate Orrery with two replication studies, each exercising a different subsystem end-to-end: (1) finding linear directions visually in the embedding space and validating them through probe training in the collection workflow (\S\ref{sec:glasgow}) and (2) activation caching and steering in the interpretability backend (\S\ref{sec:refusal}).

\subsection{Psycholinguistic Norms are Linear Directions}
\label{sec:glasgow}

LLMs self-organize the internal representation space to encode semantic dimensions linearly\cite{Marks2023TheGO}. A simple \textit{mass mean probe} can recover a factual true/false direction using multiple models $p_{\mathrm{mm}}(x)=\sigma\!\big((\mu^{+}-\mu^{-})^{\top}x\big)$ where $\mu^{\pm}$ are the class means. The pervasiveness of this linear encoding phenomenon and whether this extends to the representation spaces of sentence-embedding models  remains open. 

We probe two sets of human psycholinguistic norms in the embedding space of Gemini-Embedding-2, and MiniLM from SentenceTransformers \cite{reimers-gurevych-2019-sentence}. \textbf{Concreteness} ratings for about 40k English words and two-word expressions \cite{Brysbaert2014-py}. \textbf{Glasgow norms} ${\sim}5.5$k words rated on nine dimensions (Appendix~\ref{app:probes}).

%\begin{enumerate}
%    \item \textbf{Concreteness} ratings for about 40k English words and two-word expressions obtained from 4,000 participants \cite{Brysbaert2014-py}.
%    \item \textbf{Glasgow norms} ${\sim}5.5$k words rated on nine dimensions (arousal, valence, dominance concreteness, imageability, familiarity, age of acquisition, semantic size, and gender association).
%\end{enumerate}

The collections are first embedded, projected and colored using diverging color scales to find which dimensions appears to be linearly separated (Figure~\ref{fig:concreteness}). We then train Mass-Mean, Ridge Regression and MLP probes on the datasets, using a 80/20 split. Every dimension appearing as visually separable appears to be linearly recoverable, with concreteness TODO ADD NUMBERS 

Per-dimension results and the full protocol are in Appendix~\ref{app:probes}.


\subsection{Refusal is a Single Vector Away}
\label{sec:refusal}

%%% TODO citations:
To validate the steering backend (\S\ref{sec:sae}), we replicate
\citet{arditi2024refusallanguagemodelsmediated} on Gemma-3-4b-it (the reference uses Gemma-2). Candidate refusal directions $r_\ell$ are extracted by difference-in-means between residual-stream activations over $n{=}128$ harmful prompts from AdvBench AdvBench\cite{zou2023universaltransferableadversarialattacks} and MaliciousInstruct \cite{ICLR2024_3af25aa3} ($H^{+}$) and harmless instructions from Alpaca \cite{taori2023alpaca} ($H^{-}$), at each layer and each of the five end-of-instruction template tokens.

 Arditi removes refusal by \emph{directional ablation}, projecting $\hat r$ out of the residual stream at every layer, $h \leftarrow h - (\hat r^{\top}h)\,\hat r$. Since Gemma-3 is a pre-norm model, the residual stream grows with the number of layers \cite{hex2023residual} and this intervention, in our evaluation, collapses the model to degenerate token loops. 



%The candidate refusal directions are extracted using difference-in-mean for the residual streams over $n=128$ harmful ($H^{+}$) prompts from AdvBench\cite{zou2023universaltransferableadversarialattacks} and Malicious Instruct \cite{ICLR2024_3af25aa3} and harmless instructions from Alpaca \cite{taori2023alpaca} ($H^{-}$):
%
%We retrieve candidate vectors at each layer for each of the 5 end-of-instructions template tokens.

%\begin{equation}
%r_\ell = \tfrac{1}{|H^{+}|}\!\sum_{x\in H^{+}}\! h_\ell(x)
%       - \tfrac{1}{|H^{-}|}\!\sum_{x\in H^{-}}\! h_\ell(x).
%\label{eq:diffmeans}
%\end{equation}

%In the reference paper \emph{activation addition} was used only to induce refusal
We instead use ActAdd, the same additive mechanism as our SAE steering (\S\ref{sec:sae}), adding the unnormalized $r_{\ell^{*}}$ at a single layer: $h_{\ell^{*}} \leftarrow h_{\ell^{*}} + \alpha\,r_{\ell^{*}}$, with $\alpha=-1$ suppressing refusal and $\alpha=+1$ inducing it. At
$\ell^{*}=14$, ActAdd$(+1)$ induces refusal on \textbf{100/100} entirely benign prompts. Using Llama-Guard-4-12B \cite{inan2023llamaguardllmbasedinputoutput} as harmfulness classifier attack success rises from 0.11 to \textbf{0.83} on JailbreakBench \cite{chao2024jailbreakbench} and from 0.21 to \textbf{0.86} on HarmBench-test \cite{mazeika2024harmbenchstandardizedevaluationframework}. Completions were manually reviewed for coherence. To remove refusal a single vector on a single layer seems to be still sufficient. Since that same vector can induce refusal, classifying the input or monitoring with SAE \cite{deng2026qwenscopeturningsparsefeatures} and steering the model mid-generation seems a promising avenue to increase safety. Vector selection details, including KL-divergence on neutral prompts logits, and per-category
results are in Appendix~\ref{app:refusal}. We release the experiment code to support safety research.


%We instead use ActAdd, the same additive mechanism as our SAE steering (\S\ref{sec:sae}), which was used in the reference only to induce refusal. We use the raw norm of the residual stream vector, adding at the selected layer $\ell^{*}$, with $\alpha=-1$ to suppress refusal and $\alpha=+1$ to induce it: 

%\begin{equation}
%    h_{\ell^{*}} \;\leftarrow\; h_{\ell^{*}} + \alpha\, r_{\ell^{*}},
%    \label{eq:actadd}
%\end{equation}

%The best vector is selected by generating completions on 32 examples, evaluating induce and refusal scores by substring matching. K-L divergence is computed over the last token logits on neutral prompts to avoid approving degenerate outputs as bypassing refusal. Selecting $\ell^{*}=14$, ActAdd$(-1)$ drives the substring refusal rate to \textbf{0.00} on both JailbreakBench (from 0.84) and HarmBench-test (from 0.75), while ActAdd$(+1)$ induces refusal on \textbf{100/100} entirely benign prompts (from 4/100). Following the reference, we pair the string-match judge with a harmfulness classifier (Llama-Guard-4-12B): attack success rises from 0.11 to \textbf{0.83} on JailbreakBench and from 0.21 to \textbf{0.86} on HarmBench-test. All the response were manually reviewed and confirmed as being coherent. Detailed and per-category results are in ~\ref{app:refusal}). 

%% TODO 
%The one weak cell is disinformation (ASR 0.40/0.74, Appendix~\ref{app:refusal}): refusal still drops to zero, but the Llama-Guard under-flags persuasive-but-false text.
\section{Conclusions}

Orrery connects embedding visualization and interpretability: corpus exploration gains SAE-feature search and in-interface probing; interpretability gains the affordances of modern visualization platforms, with feature spaces rendered as navigable
collections wired to inspection and causal steering. The case studies display how the platform can enable further research: linear structure spotted visually and confirmed by probes (\S\ref{sec:glasgow}) and a refusal direction usable to test LLM-safety extracted, injected, and evaluated without leaving the platform (\S\ref{sec:refusal}).

Beyond interpretability research, the coloring, temporal, and filtering tools target the navigation of digital humanities corpora: by breaking the linear order of texts and drawing similar documents together in space, connections across authors, periods, and passages surface visually. The nebula shaders, highlight traces, and 3D labels are not ornaments but navigation aids, inviting users to treat a corpus not as a flat canvas but as a place to explore with curiosity and a sense of wonder.


%Orrery aims to improve data exploration through interpretability techniques, and to enable interpretability research offering the tooling of embedding-visualization platforms. 


%Orrery goal is twofold: bringing interpretability tooling to improve the navigation on embedding platform via SAE search, and to 


%furthermore, by allowing frictionless movement from textual collection to visualization and probing, it allows novel discovery over directions in embedding spaces. 
%Orrery aims to be both a visual experiment lab for interpretability researchers and a visual exploration canvas for people interested in the content of the textual corpus displayed. The advanced coloring, temporal analysis and filtering options are aimed at the navigation of digital humanities corpus. By breaking down the linear dimension of the books and by bringing together similar documents in space, connections across authors, periods, or quotes can be easily discovered. The nebula shaders, the highlight connected traces and the 3D topic labels are not ornaments but navigation aids, aimed at inviting users to treat datasets not as flat, anodyne canvas but place to explore with curiosity and a sense of wonder. 

%serve this same end: they are not ornament but navigation aids, inviting users to treat a dataset as a place to explore with curiosity rather than a flat, anodyne canvas. 


%embedding visualization (SAE search) and to bring the capabilities of the current embedding visualization platform (navigation, topic modeling) to interpretability research. 


%We present Orrery, an application that combines modern visualization embedding-visualisation with interpretability features and dashboards. For dataset exploration, such as the case with digital humanities corpus, it turns a textual collection into a searchable and navigable 3D constellation with topic annotation and filtrable by time period. Appendix~\ref{app:lacan} maps the complete Lacan's corpus (150k sentences). For interpretability researchers, it allows a singular end-to-end tool to solve what was previously done by several plus code: spotting linear structure visually and confirming it
%with probes (\S\ref{sec:glasgow}), or finding interesting SAE features in the galaxy, inspecting it in the dashboard, and steering the model with it
%(\S\ref{sec:refusal}). Filtering, colouring, the nebula shaders, and animation systems serve this same end: they are not ornament but navigation aids, inviting users to treat a dataset as a place to explore with curiosity rather than a flat, anodyne canvas. 


%filtering, colour, shader, and animation systems serve this same end:
%they are not ornament but navigation aids, inviting users to treat a
%dataset as a place to explore with curiosity rather than a flat,



%Future work includes exposing Qwen-Scope end-to-end
%and integrating our autointerpreter pipeline for labelling arbitrary
%SAEs.


%the pipeline used by the embedding visualisation platform and the interpretability. The platform aims at helping users explore vector dataset, in particular the digital humanities corpus. In the appendix we show the entire Lacan sentences corpus (150k sentences) and to accellerate interpretability research. For researchers orrery offers the opportunity to directly find visually linear directions in datasets and train probes from the interface, or to use the scatterplot to find interesting features and directly examine them in the dashboard and to use them to iterate in the model. The elaborate filtering, colouring, nebula shaders and animation system is not put in to be purely cosmetic: its purpose is to allow users to better navigate the space, treating the dataset like a place to explore and examine with curiousity rather than a flat anodyne canva.  


\section*{Limitations}


%Orrery's interface currently supports Gemma Scope end-to-end, while
%Qwen-Scope loads and runs at the backend level only; adding a model
%family requires porting the activation caches to its reference
%implementation, the deliberate cost of the serving-oriented backend
%(\S\ref{sec:sae}). Feature labels are fetched from Neuronpedia, so
%labelled coverage is bounded by what Neuronpedia hosts --- the
%autointerpreter pipeline we release can label arbitrary SAEs, but
%requires further validation and is not yet integrated in the frontend.
%The labels themselves are often generic or noisy, and SAE feature
%search inherits their weaknesses: although feature-based retrieval is
%motivated in the literature
%\cite{jiang2025interpretableembeddingssparseautoencoders}, we regard
%it as an experimental capability pending systematic evaluation.
%Finally, Orrery does not target the millions of points that Nomic
%Atlas or Embedding Atlas render --- a deliberate trade-off in favour
%of 3D navigation, real-time overlays, and volumetric shaders. We have
%run collections of up to 1M points on a 16\,GB machine\todo{state
%what degrades at 1M}, which covers most research corpora and every
%published SAE, but for larger datasets, 2D scale-first platforms
%remain the better tool.

Orrery at the moment only fully supports SAEs from Gemma-Scope-2 and Qwen-Scope. SAE search is also enabled only the Gemma-Scope-2 family, since the currenty application downloads labels from Neuronpedia. While the backend hosts currently an autointepreter pipeline, it has not been validated or used at scale to label Qwen features. SAE searching is also limited by the quality of the labels: autointerpreter labels are often noisy \cite{jiang2025interpretableembeddingssparseautoencoders} and could be corrected by adding post-hoc activations labels \cite{arad-etal-2025-saes}. Finally, Orrery has a scale limitation compared to the millions of points that Nomic Atlas or Embedding Atlas render. This was a deliberate design trade-off in favour of 3D navigation, real-time overlay and volumetric shaders. For larger datasets, 2D-first platforms may be preferable. 



%in the backend. This limitation is due to the lack of Neuronpedia labels for SAE and open-source libraries for SAE inference and steering. Together with the application we release a full autointepreter pipeline, which requires further validation and is not incorporated in the frontend. 

%Labels are obtained at the moment by Neuronpedia, and are generally very weak. While SAE search is motivated in the literature, it is an experimental feature that requires more 

%Topic extraction re-uses the BERTopic pipeline. While this is an improvement over Embedding Atlas density clustering plus cTF-IDF, the innovation i

%Compared to other embedding visualisation platforms, Orrery has a scale limitation of not being able to visualise millions of points at once. This was a deliberate design decision, prioritising navigation in 3D, shaders to indicate clusters, overlay traces for the highlighted points and WebGL animations. Orrery was tested up to 1 million points on a Mac Mini with 16 GB ram, which should suffice for the majority of the dataset or SAE, but for larger datasets, other platforms may be more suitable. 

%% PUT BACK LATER: final version has 1 more page
%\section*{Acknowledgments}

%I wish to thank Roksana Geworek for all the conversations, feedbacks, and inspiration on working on the geometry of the embedding spaces, Leonardo Ascione for all the help on the technical side of web-development, in particular in debugging React, and Matteo Rinaldi for wanting to bring forth developing source interpretability-libraries. 



\section*{Ethics Statement}
Our refusal case study (\S\ref{sec:refusal}) replicates a published method \cite{arditi2024refusallanguagemodelsmediated} which can be already directly executed using the refusal\_direction code available on Github\footnote{\url{https://github.com/andyrdt/refusal_direction}}. Our experiment reproduces it using Gemma-3 instead of Gemma-2, using the fork of Gemma-pytorch instead of using transformers. Our contribution is reproducibility and tooling for inspecting the safety-relevant direction. We hope that this experiment will bring further awareness that this notorious vulnerability still has not been fixed. The SAE tooling released can be used for safety research for an emerging paradigm of safety measure that rely on classification and intervention via SAE or steering vectors rather than  post-hoc alignment that can be easily broken \cite{lu2026assistantaxissituatingstabilizing}. The released code contains the extraction and evaluation pipeline, a downloaded script to retrieve the prompts, and no harmful completions. The psycholinguistic norms (\S\ref{sec:glasgow}) are public datasets of aggregate human ratings, used under their licenses.


%Our hope is to signal the need for a The psycholinguistic norms (\S\ref{sec:glasgow}) are public datasets of aggregate human ratings, used under their licenses.


%on open-weight models. It introduces no new attack capability: the technique and implementations have been public since 2024, and it can be directly reproduced using Transformers from the paper's repository. Our experiments reproduces it on Gemma-3, noting that the vulnerability still has not been fixed. 

%it only employs the fork of Gemma-Pytorch instead of using the Transformers implementation. Our contribution is reproducibility and tooling for \emph{inspecting} safety-relevant directions. The released experiment code contains the extraction and evaluation pipeline, a downloaded to retrieve the prompts from the repository of Arditi, and no harmful completions. Our hope is to signal the need for a paradigm of safety measure that rely on classification and intervention rather than easily broken post-hoc alignemnt. 


%and it applies
%only to models whose weights the user already fully controls, so it
%does not transfer to models served behind an API. Our contribution is
%reproducibility and tooling for \emph{inspecting} safety-relevant
%directions --- the workflow that motivates the monitoring-and-steering
%argument in \S\ref{sec:refusal}. The released experiment code contains
%the extraction and evaluation pipeline, prompts, and judge scores, but
%no harmful completions.\todo{confirm this is true of the repo before
%claiming it} All steered outputs were manually reviewed; we note that
%Llama-Guard under-flags disinformation, and report it as a judge
%limitation. The psycholinguistic norms (\S\ref{sec:glasgow}) are
%public datasets of aggregate human ratings, used under their licenses.





\bibliography{custom}

\appendix

\section{Embedding Settings}
\label{app:embedding}

The platform supports 8 diffent providers for embeddings. Local: SentenceTransformers (default), Ollama, FlagEmbedding for BGE~\cite{chen-etal-2024-m3}, and PyTorch for Qwen~\cite{zhang2025qwen3embeddingadvancingtext}; via API key: OpenAI, Cohere, Gemini, and HuggingFace. The settings include task-type prompts for models that use them, such as EmbeddingGemma~\cite{vera2025embeddinggemmapowerfullightweighttext}. 


\section{Refusal-direction evaluation}
\label{app:refusal}

The selected direction is \texttt{pre\_attn}, token position $-3$, layer 14, chosen from the difference-in-means candidates of Eq.~\ref{eq:diffmeans} (128 harmful / 128 harmless instructions, bare Gemma chat template, no system prompt). Behavioural evaluation greedy-generates 256-token completions and scores each with (i)~the reference substring refusal judge and (ii)~Llama-Guard-4-12b attack-success rate (ASR). We omit the ablation condition from harm scoring: its 0\% refusal is an artifact of the collapse above (degenerate loops the substring judge miscounts as non-refusal), which is precisely why the safety classifier is reported alongside.


\begin{table}[t]
\centering
{\small
\setlength{\tabcolsep}{4pt}
\begin{tabular}{@{}llccc@{}}
\toprule
Set ($n$) &      & base & ActAdd$_{-1}$ & ActAdd$_{+1}$ \\
\midrule
JBB (100)      & ref. & 0.84 & \textbf{0.00} & ---           \\
               & ASR  & 0.11 & \textbf{0.83} & ---           \\
HB-test (159)  & ref. & 0.75 & \textbf{0.00} & ---           \\
               & ASR  & 0.21 & \textbf{0.86} & ---           \\
Harmless (100) & ref. & 0.04 & ---           & \textbf{1.00} \\
\bottomrule
\end{tabular}}
\caption{Refusal rate (ref.) and harm attack-success rate (ASR) under activation addition. ActAdd$_{-1}$ suppresses refusal, ActAdd$_{+1}$ induces it; a single direction controls both. JBB: JailbreakBench; HB-test: HarmBench-test.}
\label{tab:refusal-main}
\end{table}

\begin{table}[t]
\centering
{\small
\begin{tabular}{@{}lccc@{}}
\toprule
HarmBench category & $n$ & base & ActAdd$_{-1}$ \\
\midrule
Chem./bio.      & 19 & 0.37 & \textbf{1.00} \\
Cybercrime      & 33 & 0.45 & \textbf{1.00} \\
Harassment      & 16 & 0.00 & 0.88 \\
Illegal         & 47 & 0.15 & 0.87 \\
Misinfo         & 27 & 0.15 & 0.74 \\
Harmful (gen.)  & 17 & 0.00 & 0.59 \\
\bottomrule
\end{tabular}}
\caption{Per-category harm ASR (HarmBench-test) at $\ell^{*}=14$. The weakest cells are disinformation and generic ``harmful'', where the model complies but Llama-Guard under-scores the output.}
\label{tab:refusal-cat}
\end{table}






\section{Probe Settings and Per-Dimension Results}
\label{app:probes}

\paragraph{Probe settings.}
All probes share one $80/20$ split (seed $7$) on per-dimension $z$-scored embeddings. \textsc{Ridge}: $\ell_2$ linear regression, $\alpha=1$.
\textsc{SVR}: support-vector regression, RBF kernel, $C=1$. \textsc{MLP}: one
hidden layer ($256$, ReLU, dropout $0.1$), Adam (lr $10^{-3}$, weight decay
$10^{-4}$), MSE loss, batch size $64$, up to $100$ epochs with patience-$10$
early stopping on validation loss. \textsc{Mass-mean}: the closed-form
unit-normalised difference-of-means direction $\boldsymbol{\mu}_{+}-\boldsymbol{\mu}_{-}$ of
\citet{Marks2023TheGO} on a median split, scored by projection (reported as
Spearman $\rho$, an uncalibrated ranking). \textsc{Logistic}: $\ell_2$ logistic
regression ($C=1$, up to $1000$ iterations) on a median split of the norm,
reported as classification accuracy. Ridge/SVR/MLP report $R^2$. All figures are
the \emph{last} hidden layer (MiniLM L6, EmbeddingGemma L24) --- i.e.\ the model
output embedding. \textsc{Concreteness-40k} is the full Brysbaert set
($39{,}954$ words) on the final pooled MiniLM embedding.


\begin{figure*}[t]
    \centering
    \includegraphics[width=\textwidth]{media/EMNLP.png}
    \caption{pipeline}
    \label{fig:emnlp}
\end{figure*}


\begin{table*}[t]
\centering
{\small
\begin{tabular}{@{}lccccc@{}}
\toprule
\textbf{Norm} & \textbf{Ridge} & \textbf{SVR} & \textbf{MLP}
  & \textbf{Mass-mean} & \textbf{Logistic} \\
              & $R^2$ & $R^2$ & $R^2$ & $\rho$ & acc \\
\midrule
\multicolumn{6}{@{}l}{\emph{MiniLM-384}} \\
concreteness     & 0.68 & 0.71 & 0.71 & 0.70 & 0.85 \\
imageability     & 0.63 & 0.66 & 0.66 & 0.67 & 0.80 \\
valence          & 0.59 & 0.68 & 0.66 & 0.67 & 0.79 \\
aoa              & 0.56 & 0.60 & 0.58 & 0.57 & 0.78 \\
semsize          & 0.53 & 0.60 & 0.57 & 0.62 & 0.79 \\
gender           & 0.49 & 0.59 & 0.53 & 0.65 & 0.76 \\
familiarity      & 0.51 & 0.50 & 0.50 & 0.60 & 0.79 \\
arousal          & 0.43 & 0.51 & 0.46 & 0.54 & 0.74 \\
dominance        & 0.40 & 0.50 & 0.45 & 0.53 & 0.71 \\
concreteness-40k & 0.70 & 0.72 & 0.74 & 0.80 & 0.84 \\
\midrule
\multicolumn{6}{@{}l}{\emph{EmbeddingGemma-768}} \\
concreteness     & 0.67 & 0.71 & 0.61 & 0.80 & 0.81 \\
imageability     & 0.62 & 0.64 & 0.55 & 0.76 & 0.80 \\
valence          & 0.60 & 0.68 & 0.57 & 0.77 & 0.80 \\
aoa              & 0.54 & 0.53 & 0.34 & 0.66 & 0.78 \\
semsize          & 0.57 & 0.61 & 0.48 & 0.71 & 0.78 \\
gender           & 0.49 & 0.61 & 0.35 & 0.65 & 0.75 \\
familiarity      & 0.42 & 0.49 & 0.14 & 0.65 & 0.75 \\
arousal          & 0.43 & 0.53 & 0.31 & 0.61 & 0.72 \\
dominance        & 0.45 & 0.52 & 0.32 & 0.63 & 0.71 \\
\bottomrule
\end{tabular}}
\caption{Per-dimension recovery of the psycholinguistic norms at the last hidden layer (output embedding), validation split. Ridge/SVR/MLP report $R^2$; mass-mean is Spearman $\rho$ (an uncalibrated ranking, not comparable to $R^2$); logistic is median-split accuracy. Every norm is recovered well above chance by the linear Ridge probe, with SVR and the MLP adding little --- consistent with a largely linear encoding. Concreteness is the most strongly recovered dimension and the affective norms (arousal, dominance) the least; rows are roughly ordered by encodability.}
\label{tab:probes}
\end{table*}








\end{document}



%\begin{figure*}[!t]
%\centering
%\begin{subfigure}[b]{0.48\textwidth}
%\centering
%\includegraphics[width=\textwidth]{media/safeunsafewhite.png}
%\caption{Safety Label}
%\label{fig:safety_fv}
%\end{subfigure}
%\hfill
%\begin{subfigure}[b]{0.48\textwidth}
%\centering
%\includegraphics[width=\textwidth]{media/concretenesswhite.png}
%\caption{Literary Format}
%\label{fig:format_fv}
%\end{subfigure}
%\caption{lorem te ipsum}
%\label{fig:clustering_fv}
%\end{figure*}